\section{Principes}

\begin{frame}
\frametitle{Principe du partitionnement par hachage}

On prend la clé; on l'interprète comme un entier, on calcule le
reste $r$ de la division par $N$: le document est affecté au fragment $r$.

\vfill

\figSlideWithSize{partition-hachage}{10cm}

\vfill

\begin{remblock}{}Simple mais...
Fonctionne très bien pour $N$, le nombre de fragments, \alert{fixe}.
\end{remblock}
\end{frame}



\begin{frame}
\frametitle{Les opérations}

Implantation très simple et très efficace.

\begin{redbullet}
  \item $get(i)$: calculer $r = h(i)$, et accéder au  
    fragment dont l'adresse est $a_r$, chercher le document en mémoire;
  \item $put(i, d)$: calculer $r = h(i)$, 
    insérer $d$ dans le fragment dont l'adresse est  $a_r$;
  \item $delete(i)$: comme la recherche, avec effacement du document trouvé;
  \item $range(i, j)$: \alert{pas possible avec une structure par hachage}, il 
    faut faire un parcours séquentiel complet.
\end{redbullet}     

\vfill

\begin{remblock}{Elasticité}
Si $N$ change: tout devient faux. 
\end{remblock}
\end{frame}


\begin{frame}
\frametitle{Le problème}


Tout repose  sur l'hypothèse que la fonction $h()$ est \alert{immuable}.

\vfill
 Soit $N=5$, donc  $h(i) = i \mod 5$. Soit un flux
continu d'insertion et de recherche de documents, parmi lesquelles
l'insertion, puis la recherche de l'identifiant 17. 
 
\begin{redbullet}
   \item  on effectue \textit{put(17, d)}, $d$ est inséré dans $F_2$ (tout le monde suit?) ;
  \item les insertions continuent, jusqu'à la nécessité d'ajouter
      un sixième fragment: maintenant $h(i) = i \mod 6$.
   \item on effectue $get(17)$, on cherche dans $F_5$
      (vous voyez pourquoi?) qui ne contient pas $d$.
\end{redbullet}      

\alert{Pas de solution simple!!}
\end{frame}

\begin{frame}
\frametitle{Systèmes distribués basés sur le hachage}

Idée simple\,: tout le monde utilise le même hachage, les ``blocs''
sont les serveurs.

\vfill

Le système est nécessairement \alert{dynamique} (ajout  de serveurs), donc
on ne peut pas utiliser une fonction de hachage basée directement sur $N$, 
qui varie.


\vfill

En 1997, un article a proposé une solution maintenant très largement utilisé\,: le
\alert{hachage cohérent} ou \textit{Consistent haching}.


\begin{redbullet}
  \item D'abord conçu et utilisé pour des \alert{caches distribués} (\textit{memcached})

  \item Popularisé pour la gestion de données par le système Dynamo (Amazon, 2007)

  \item Maintenant intégré à de très nombreux systèmes\,: Voldemort, Riak, Cassandra, ...
 
\end{redbullet}
\end{frame}

\section{Le hachage cohérent}

\begin{frame}
\frametitle{Hachage cohérent (\textit{consistent hashing})}

On  prend une valeur de $N$ très grande et immuable, p.e. $2^{64}$

\begin{redbullet}
  \item la fonction de hachage  $h(v) = v \mod N $ est \alert{fixe},
    appliquée aux clés \alert{et} aux serveurs (identifiées par leur IP)
          vers l'espace d'adressage  $A=[0, 2^{64}-1$];
  \item $A$ est organisé comme un anneau parcouru dans le sens des aiguilles d'une montre\,;
\end{redbullet}

\vfill

Serveurs et document sont un emplacement dans $A$ défini par $h$.
\alert{Il n'y a \alert{aucune chance} qu'un serveur ait la même valeur de hachage qu'un 
serveur, d'où}
\vfill

\begin{remblock}{Régle d'affectation des documents aux serveurs}
Si $S$ et $S'$ sont deux serveurs adjacents sur l'anneau, toutes les clés de l'intervalle 
   $]h(S), h(S')]$ appartiennent à $S'$.
\end{remblock}
\end{frame}

\begin{frame}
\frametitle{Illustration}

Serveurs \alert{et} documents sont placés sur l'anneau.

\vfill

\figSlideWithSize{ch-general}{8cm}

\vfill

Notez le déséquilibre des fragments (pour l'instant).
\end{frame}


\begin{frame}
\frametitle{La table de hachage}

% A two-columns presentation
\begin{tabular}{l|p{6cm}}
   \parbox{5cm}{
   
    Elle établit une correspondance
entre le découpage de l'anneau en arcs de cercle, et l'association de chaque
arc à un serveur.

\medskip

\begin{tabular}{r|l}
Intervalle  & Serveur\\
\hline
$]c, a]$ & S1 \\
$]a, b]$ & S2  \\
$]b, c]$ & S3  \\
\hline
\end{tabular}

   }  % Inclusion of an XML file
  & 
  \parbox{6.5cm}{
   \parbox{5cm}{\figSlideWithSize{ch-general}{6cm} }
}

\end{tabular}

\end{frame}


\begin{frame}
\frametitle{Quelques détails vraiment utiles}

Comment fonctionne la tolérance aux pannes\,? Comment fonctionne l'équilibrage\,?


\vfill

% A two-columns presentation
\begin{tabular}{l|p{6cm}}
   \parbox{6cm}{


\alert{Panne} $\Rightarrow$ géré par la réplication (surprise!)\,; 
par exemple on copie sur la machine suivante dans l'anneau, etc.


  \vfill
  
\alert{Equilibrage} $\Rightarrow$ un même serveur (physique) est distribué
en plusieurs points (virtuels) sur l'anneau.

\begin{redbullet}
  \item plus il y a de points, plus le serveur recevra de données\,;
  \item en cas de panne, réoganisation répartie plus justement.
\end{redbullet}
  
   }  % Inclusion of an XML file
  & 
  \parbox{6.5cm}{\figSlideWithSize{ch-multipos}{6cm} }
\end{tabular}

\end{frame}


\begin{frame}
\frametitle{Ajout et suppression de serveurs}


\figSlideWithSize{ch-ajout}{12cm}

\vfill

Une réorganisation \alert{locale} est suffisante ($S_3$ transmet à $S_4$).

\end{frame}

\section{Elastic Search}


\begin{frame}
\frametitle{Mais où est donc la table de hachage\,?}

\alert{Qui nous dit quelle est la liste des serveurs et leur intervalle}\,?

\smallskip

Il nous faut un \alert{routeur}.

\smallskip

\figSlideWithSize{archi-consistenthash}{8cm} 
\smallskip

Autres solutions\,: (i) chaque nœud est un routeur (acceptable
quand on n'ajoute/retire pas de serveurs tout le temps), (ii) le routeur est intégré au client.

\end{frame}


